% ============================================================
% Section 9 — Where Should Evidence Signals Enter the Agent? (draft v1, 2026-08-31)
% 数据: HANDOFF §6(Judge-C, 选择性预测, best-of-4)+ §9(J8 注入消融)
% 复核(2026-08-31 本窗口): rerank_results2.json = 237 题 × 4 rollouts,
%   majority 28.3 / oracle 35.0 复算吻合。
% 措辞纪律: J8 用 "no reliable improvement, trend negative"。
% ============================================================

% 2026-09-20: 整节移入附录(用户决议: 安置研究不作为贡献)
\section{Where should evidence signals enter the agent?}
\label{sec:placement}

The evidence signal of Eq.~\ref{eq:coverage} enters training at the
query step (\S\ref{sec:method-credit}) or in the trajectory return
(\S\ref{sec:method-effect}); both are one placement of a more general
resource. Evidence supervision can enter the agent at four
points: as a training reward, as an offline audit of finished
trajectories, as a runtime controller that selects among candidate
rollouts, and as an observation written into the agent's context. We
tested all four. Gold annotations exist only at training time, so the
placements that run at inference use a 4B verifier distilled from the
same annotations (held-out AUC 0.996, zero-shot 2Wiki 0.987; component
validation in Appendix~\ref{app:verifier}); on full policy rollouts it predicts answer
correctness at AUC 0.742.

\paragraph{Training reward: works.} As a reward the signal produces
the deep-hop gains of \S\ref{sec:results}, through the query-step
credit of \S\ref{sec:method-credit} and, more weakly, through the
trajectory ranking of \S\ref{sec:method-effect}
(\S\ref{sec:mechanism}).

\paragraph{Offline audit: works.} Scoring finished trajectories,
evidence coverage supports selective
prediction~\citep{geifman2017selective}: on a stratified audit set of
600 trajectories with 49\% base accuracy, answering only the 5\%
ranked most-covered yields 79\% accuracy, and the 10\% level yields
63\%. The subset whose plan is fully resolved---5.8\% of
questions---is 82\% correct.

\paragraph{Runtime control: fails.} Reranking four sampled rollouts
per question by verifier-scored evidence coverage yields 26.6--27.0
EM across every scoring strategy on a 237-question audit
set---below majority vote at 28.3, with the best-of-4 oracle at
35.0. The trajectories explain the failure: even correct rollouts
fill their declared plan slots in order at rates of only
60/19/7.5/2.4\% for slots one through four, so plan-anchored
evidence scoring systematically undervalues rollouts that reach the
answer off-plan.

\paragraph{Injected observation: fails.} From a common checkpoint,
two variants trained in the same window under identical guard criteria:
one receives a verifier-computed slot-status line in every
observation, one does not. At the last common evaluation step the
injected variant scores 46.71 versus 47.83 for the control on the full
dev set ($z=-1.80$; 99:126): no reliable improvement, and the trend
is negative. The dynamics show why. Injection costs $-4.8$ points at
twenty episodes, and the policy then adapts back to parity---it
learns to ignore the extra line rather than exploit it. The budget
line already present in the observation carries the
stopping-relevant information, and the verifier's conservative bias
(the plan--path misalignment above) leaves little marginal signal to
use.

\paragraph{Reading.} The pattern follows from what the signal
measures. Coverage separates correct from incorrect rollouts within a
group at AUC 0.66--0.73 (\S\ref{sec:mech-not}); at inference only a
verifier's estimate of it is available. Progress supervision pays where
% 原句: "Coverage quantifies evidence progress, not answer correctness
% (\S\ref{sec:mech-not})."
it reorders learning or stratifies risk after the fact; it has no
purchase where a single decision must be steered at inference, and
the policy's divergence between declared plans and executed paths
degrades every plan-anchored use. Evidence supervision is
placement-sensitive: the same deterministic signal, unchanged, moves
from decisive to useless as it moves across these four entry points.
